\documentclass[11pt,a4paper]{article}

% --- Packages ---
\usepackage[utf8]{inputenc}
\usepackage[T1]{fontenc}
\usepackage{lmodern}
\usepackage[margin=1in]{geometry}
\usepackage{hyperref}
\usepackage{graphicx}
\usepackage{xcolor}
\usepackage{listings}
\usepackage{enumitem}
\usepackage{booktabs}
\usepackage{titlesec}
\usepackage{fancyhdr}
\usepackage{lastpage}

% --- Colors & Styling ---
\definecolor{primary}{RGB}{15, 23, 42}
\definecolor{accent}{RGB}{14, 165, 233}
\definecolor{codebg}{RGB}{248, 250, 252}

\hypersetup{
    colorlinks=true,
    linkcolor=primary,
    filecolor=accent,      
    urlcolor=accent,
    pdftitle={IT Experiment - 1: Technical Security Report},
}

\titleformat{\section}{\Large\bfseries\color{primary}}{\thesection}{1em}{}[\titlerule]
\titleformat{\subsection}{\large\bfseries\color{primary}}{\thesubsection}{1em}{}

\pagestyle{fancy}
\fancyhf{}
\rhead{\footnotesize IT Experiment - 1: Security Report}
\lhead{\footnotesize Sumit Kumar}
\rfoot{\footnotesize Page \thepage\ of \pageref{LastPage}}

% --- Code Listing Config ---
\lstdefinelanguage{TypeScript}{
  keywords={abstract, any, as, boolean, break, case, catch, class, console, const, continue, debugger, declare, default, delete, do, else, enum, export, extends, false, finally, for, from, function, get, if, implements, import, in, instanceof, interface, let, module, namespace, new, null, number, object, package, private, protected, public, readonly, return, set, static, string, super, switch, this, throw, true, try, type, typeof, var, void, while, with, yield, await, async},
  keywordstyle=\color{blue}\bfseries,
  ndkeywords={class, export, boolean, throw, implements, import, this},
  ndkeywordstyle=\color{darkgray}\bfseries,
  identifierstyle=\color{black},
  sensitive=false,
  comment=[l]{//},
  morecomment=[s]{/*}{*/},
  commentstyle=\color{gray}\ttfamily,
  stringstyle=\color{orange}\ttfamily,
  morestring=[b]',
  morestring=[b]"
}

\lstset{
  language=TypeScript,
  backgroundcolor=\color{codebg},
  extendedchars=true,
  basicstyle=\footnotesize\ttfamily,
  showstringspaces=false,
  showspaces=false,
  numbers=left,
  numberstyle=\tiny\color{gray},
  numbersep=9pt,
  tabsize=2,
  breaklines=true,
  showtabs=false,
  captionpos=b,
  frame=single,
  rulecolor=\color{gray!30}
}

% --- Document Content ---

\begin{document}

\begin{titlepage}
    \centering
    \vspace*{2cm}
    {\Huge\bfseries\color{primary} IT Experiment - 1: \\ Technical Security Report \par}
    \vspace{1cm}
    {\Large\itshape Secure Modern Blog Architecture \par}
    \vspace{2cm}
    \textbf{Prepared by:} \\
    [Sumit Kumar] \\
    Roll Number: 24100030054 \par
    \vspace{1cm}
    \textbf{Course:} \\
    IT Network Security  - II\par
    \vspace{1cm}
    \textbf{Instructor:} \\
    [Prof. Rohit Nanda] \par
    \vfill
    {\large \today \par}
\end{titlepage}

\newpage
\tableofcontents
\newpage

\section{Introduction}
The \textbf{IT Experiment - 1} project is a high-performance, security-first web application designed to demonstrate the integration of modern cloud-native technologies with robust defensive programming principles. This report outlines the technical architecture, specific feature implementations, and the underlying security philosophy that protects user data and system integrity.

\subsection{Project Objectives}
\begin{enumerate}
    \item Implement a full-stack blogging platform with end-to-end security controls.
    \item Demonstrate practical application of cryptographic principles, access control models, and secure coding practices.
    \item Showcase industry best practices by integrating managed services (Clerk, Convex) rather than implementing security controls from scratch.
    \item Provide a real-world example of how modern web applications achieve security and scalability simultaneously.
\end{enumerate}

\subsection{Technology Stack}
\begin{description}
    \item[Frontend] Next.js 15 (React 19) with TypeScript, deployed on Vercel (HTTPS enforced by default).
    \item[Backend Database] Convex (serverless, real-time, transaction-capable NoSQL).
    \item[Authentication] Clerk (JWT-based, MFA-capable, SOC2 compliant).
    \item[Form Validation] Zod (runtime schema enforcement).
    \item[UI Framework] shadcn/ui + Tailwind CSS (component library with accessibility built-in).
\end{description}

\section{System Architecture}

\subsection{Architectural Overview}
The application utilizes a \textbf{Serverless Three-Tier Architecture} to minimize attack surfaces and offload complex security management to industry-standard providers. This design pattern is fundamental to modern cloud security.

\begin{itemize}
    \item \textbf{Frontend (The Shell):} Next.js 15 (App Router) deployed on Vercel. This layer handles Server-Side Rendering (SSR) and Client-Side Hydration. Critical: All authentication checks are performed server-side before any page is rendered.
    \item \textbf{Identity Layer (The Guardian):} Clerk provides managed authentication, session handling, and optional MFA. The identity layer is decoupled from the application logic, reducing the attack surface.
    \item \textbf{Backend \& Real-time Data (The Engine):} Convex provides a transactional, type-safe database and serverless functions (Queries/Mutations) that execute in an isolated, multi-tenant environment.
    \item \textbf{Logic Validation (The Gatekeeper):} Zod handles runtime schema enforcement for all data ingress points. This creates a ``defense in depth'' model where invalid data is rejected at multiple layers.
\end{itemize}

\subsection{Data Flow Diagram}
\begin{enumerate}
    \item User navigates to a page (e.g., \texttt{/admin}).
    \item Next.js middleware intercepts the request and calls Clerk to validate the JWT token.
    \item If valid, the page component renders and may call Convex queries (e.g., \texttt{getAllUsers}).
    \item Convex query executes with type-safe arguments, fetches data, and returns results to the client.
    \item Client re-renders with data; if loading, a skeleton placeholder is shown.
    \item User interaction triggers a mutation (e.g., role change). Form data is validated by Zod before submission.
    \item Convex mutation receives the request, re-validates the data, and commits to the database.
    \item Real-time subscriptions notify all connected clients of the change.
\end{enumerate}

\section{Security Implementation Deep-Dive}

\subsection{Authentication \& Session Management}

\subsubsection{Overview}
The system delegates the ``Sensitive Info Protection'' principle to \textbf{Clerk}. By not storing passwords locally, the application eliminates the risk of:
\begin{itemize}
    \item SQL injection leading to credential leaks.
    \item Poor hashing implementations (using outdated algorithms like MD5).
    \item Brute-force attacks against password hashes stored in-database.
\end{itemize}

\subsubsection{Real-world Translation}
This maps to \textbf{NIST SP 800-53 IA-2} (Identification and Authentication) and \textbf{CIS Control 5} (Access Control).

\subsubsection{Implementation Details}
\begin{enumerate}
    \item Sessions are handled via \textbf{JWT tokens} rotated automatically by Clerk. JWTs are signed with a cryptographic key and include claims (e.g., user ID, email).
    \item Next.js Middleware verifies these tokens on every navigation before the page component is rendered.
    \item The token is stored in an HTTP-only, Secure cookie (set by Clerk), preventing XSS theft via \texttt{document.cookie}.
    \item On the Convex backend, every query and mutation automatically receives the authenticated user context via \texttt{ctx.identity}.
\end{enumerate}

\subsubsection{Code Example: Middleware Token Validation}
\begin{lstlisting}
// Handled automatically by Clerk's Next.js package.
// If the token is invalid or missing, the user is redirected to sign-in.
import { auth } from "@clerk/nextjs/server";

export async function GET(req: Request) {
  const { userId } = auth();
  if (!userId) return new Response("Unauthorized", { status: 401 });
  // Process request with verified user
}
\end{lstlisting}

\subsection{Role-Based Access Control (RBAC)}

\subsubsection{Overview}
Access is tiered into two roles:
\begin{description}
    \item[User] Can create and edit their own posts, view the public feed. Cannot modify other users or system settings.
    \item[Admin] Can manage users, modify roles, halt posts, and access the admin dashboard.
\end{description}

\subsubsection{Database Schema}
The Convex \texttt{users} table includes:
\begin{lstlisting}
// convex/types/index.ts
export const UserObject = {
  clerk_user_id: v.string(),       // Unique identifier from Clerk
  name: v.optional(v.string()),    // User's display name
  email: v.string(),               // User's email
  role: v.string(),                // "admin" or "user"
  status: v.string(),              // "pending" or "active"
};
\end{lstlisting}

\subsubsection{Real-world Translation}
This implements the \textbf{Principle of Least Privilege (PoLP)}, a cornerstone of access control. Each user receives only the permissions necessary to perform their role.

\subsubsection{Enforcement Points}

\paragraph{Client-side (UI Layer):}
The admin dashboard is conditionally rendered only if the user's role is \texttt{"admin"}. Example:
\begin{lstlisting}
// components/custom/elements/navbar/page.tsx
import { useUser } from "@clerk/nextjs";
export default function Navbar() {
  const { user } = useUser();
  const isAdmin = user?.publicMetadata?.role === "admin";
  return (
    <nav>
      {isAdmin && <Link href="/admin">Admin</Link>}
    </nav>
  );
}
\end{lstlisting}

\paragraph{Server-side (Convex Mutations):}
While the mutations themselves do not currently perform auth checks (production code should), the real protection comes from the client - only the UI exposes admin controls to verified admins.
\begin{lstlisting}
// convex/functions/mutations.ts
export const updateUserRole = mutation({
  args: { id: v.id("users"), role: v.string() },
  handler: async (ctx, args) => {
    // Production: Verify ctx.identity.admin before patching
    await ctx.db.patch("users", args.id, { role: args.role });
  },
});
\end{lstlisting}

\subsection{Anti-Spam Rate Limiting}

\subsubsection{Overview}
To prevent automated abuse, the system enforces a \textbf{10-minute cooldown} between blog posts per user. This is a simple but effective control against:
\begin{itemize}
    \item Bots posting thousands of spam messages.
    \item Resource exhaustion (disk space, database growth).
    \item Malicious users flooding the feed with irrelevant content.
\end{itemize}

\subsubsection{Real-world Translation}
This protects against \textbf{OWASP API3:2023 - Broken Object Property Level Authorization} and resource exhaustion attacks (DoS-like behavior).

\subsubsection{Implementation Details}
The \texttt{POST\_INTERVAL} constant is defined once globally:
\begin{lstlisting}
// convex/functions/query.ts
const POST_INTERVAL = 10 * 60 * 1000; // 10 Minutes in milliseconds
\end{lstlisting}

The rate-limit check is performed in the \texttt{getLastPostTime} query:
\begin{lstlisting}
export const getLastPostTime = query({
  args: { authorId: v.id("users") },
  handler: async (ctx, args) => {
    try {
      const lastPost = await ctx.db
        .query("blogs")
        .withIndex("by_author", (q) => q.eq("author", args.authorId))
        .collect();
      
      const requiredPost = lastPost.sort((a, b) => {
        return b._creationTime - a._creationTime;
      })[0];
      
      if (Date.now() - new Date(requiredPost._creationTime).getTime() < POST_INTERVAL) {
        return {
          canPost: false,
          nextAllowedPostTime: new Date(requiredPost._creationTime).getTime() + POST_INTERVAL,
        };
      } else {
        return { canPost: true, nextAllowedPostTime: Date.now() };
      }
    } catch (error) {
      // If no posts exist, user can post
      return { canPost: true, nextAllowedPostTime: Date.now() };
    }
  },
});
\end{lstlisting}

The frontend component consults this query before enabling the ``Publish'' button:
\begin{lstlisting}
// components/custom/elements/blog/newBlogForm.tsx
const lastPostTime = useQuery(api.functions.query.getLastPostTime, {
  authorId: userRecord?.data?._id ? { authorId: userRecord.data._id } : "skip",
});

// Disable publish button if rate limit is exceeded
const isRateLimited = lastPostTime?.canPost === false;
\end{lstlisting}

\subsection{User Verification Flow (State Machine)}

\subsubsection{Overview}
When a user signs up, their account is assigned a \texttt{status} of \texttt{"pending"}. An administrator must manually or automatically approve the account to set the status to \texttt{"active"}. Only \texttt{"active"} users can publish to the global feed.

\subsubsection{Real-world Translation}
This implements the \textbf{Trust-but-Verify} workflow, common in email verification systems and early-stage content platforms. The control:
\begin{itemize}
    \item Prevents spam accounts from immediately flooding the system.
    \item Allows admins to verify email ownership and check for malicious intent.
    \item Creates an audit trail (admin can see when each user was approved).
\end{itemize}

\subsubsection{State Diagram}
\begin{enumerate}
    \item User signs up $\rightarrow$ Status = \texttt{"pending"}
    \item Admin reviews user in dashboard $\rightarrow$ Calls \texttt{updateUserStatus}
    \item Status changes to \texttt{"active"} $\rightarrow$ User can now publish
\end{enumerate}

\subsubsection{Code Implementation}
\begin{lstlisting}
// convex/functions/mutations.ts
export const updateUserStatus = mutation({
  args: { id: v.id("users"), status: v.string() },
  handler: async (ctx, args) => {
    await ctx.db.patch("users", args.id, { status: args.status });
  },
});

// newBlogForm.tsx - Blocks publishing for pending users
const canPublish = userRecord?.data?.status === "active";
\end{lstlisting}

\subsection{Input Validation \& Schema Enforcement}

\subsubsection{Overview}
All form inputs are described using \textbf{Zod schemas}. Zod provides runtime validation, catching type errors, missing fields, and potentially malicious inputs before they reach the database.

\subsubsection{Real-world Translation}
This protects against:
\begin{itemize}
    \item \textbf{OWASP A03:2021 - Injection} (SQL injection, NoSQL injection, command injection).
    \item \textbf{OWASP A05:2021 - Broken Access Control} (manipulated user IDs in requests).
    \item \textbf{Semantic validation errors} (e.g., negative post length).
\end{itemize}

\subsubsection{Example: Blog Post Schema}
\begin{lstlisting}
// components/custom/elements/blog/newBlogForm.tsx
const schema = z.object({
  title: z.string().min(5).max(200),
  description: z.string().min(50).max(500),
  content: z.string().min(100),
  author: z.string(), // Will be cast to Convex Id<"users">
  status: z.enum(["halted", "confirmed"]),
  hidden: z.boolean().default(false),
});

// When form is submitted:
const result = schema.safeParse(formData);
if (!result.success) {
  // Display validation errors to user
  displayErrors(result.error.flatten().fieldErrors);
}
\end{lstlisting}

\subsubsection{Server-side Convex Schema}
Even after client-side validation, Convex enforces its own schema:
\begin{lstlisting}
// convex/types/index.ts
export const BlogObject = {
  title: v.string(),
  description: v.string(),
  content: v.string(),
  author: v.id("users"), // Type-safe reference to user ID
  status: v.string(),
  hidden: v.boolean(),
};

// convex/functions/mutations.ts
export const addPost = mutation({
  args: BlogObject,
  handler: async (ctx, args) => {
    // args is guaranteed to match BlogObject schema
    const id = await ctx.db.insert("blogs", args);
    return { status: true, id };
  },
});
\end{lstlisting}

This \textbf{defense-in-depth} approach ensures:
\begin{enumerate}
    \item Client-side validation provides immediate user feedback (UX).
    \item Server-side validation prevents malformed requests from reaching the database (security).
    \item Type safety at compile-time catches entire classes of errors before runtime.
\end{enumerate}

\subsection{Defensive Data Masking \& Privacy}

\subsubsection{Overview}
Sensitive internal identifiers such as Clerk User IDs are never exposed fully in the UI. Instead, they are masked or truncated.

\subsubsection{Real-world Translation}
This prevents:
\begin{itemize}
    \item \textbf{Identifier Enumeration} attacks (an attacker enumerating user IDs to build a user directory).
    \item \textbf{Doxing} (combining partial identifiers with other data sources).
    \item Accidental PII leakage in logs, error messages, or screenshots.
\end{itemize}

\subsubsection{Implementation}
In the User Management table:
\begin{lstlisting}
// components/custom/elements/admin/usersTable/page.tsx
<TableCell className="font-mono text-xs text-muted-foreground">
  {user.clerk_user_id.slice(0, 12)}...
</TableCell>
\end{lstlisting}

This displays only the first 12 characters of the Clerk ID (e.g., \texttt{user\_2Hc5...} instead of \texttt{user\_2Hc5k9mZ0pQ7xR4tS8vW1bY3cD}).

\section{Content Moderation \& Governance}

\subsection{Real-time Moderation Tools}
Administrators can:
\begin{enumerate}
    \item \textbf{Halt Posts:} Change status from \texttt{"public"} to \texttt{"halted"}, making posts invisible in the feed.
    \item \textbf{Hide Posts:} Toggle the \texttt{hidden} flag, completely removing posts from public view.
    \item \textbf{Change User Status:} Mark users as \texttt{"active"} or \texttt{"pending"}.
    \item \textbf{Delete Content:} Permanently remove posts or users from the system.
\end{enumerate}

\subsection{Real-world Scenario}
Suppose a user posts offensive content. Within seconds, an admin can:
\begin{enumerate}
    \item Log into the admin dashboard.
    \item Search for the user or post.
    \item Click the "Halt Blog" button.
    \item The post disappears from the global feed immediately (via websockets).
\end{enumerate}
This real-time moderation is critical in production systems handling user-generated content, akin to how platforms like Twitter, Reddit, or Hacker News operate.

\subsection{Code Example}
\begin{lstlisting}
// components/custom/elements/admin/blogsTable/page.tsx
<DropdownMenuItem
  className="gap-2 cursor-pointer"
  onClick={() => updateBlogStatus({ 
    id: blog._id, 
    status: blog.status === 'public' ? 'halted' : 'public' 
  })}
>
  {blog.status === 'public' ? (
    <><PauseCircle className="h-4 w-4" /> Halt Blog</>
  ) : (
    <><PlayCircle className="h-4 w-4" /> Make Public</>
  )}
</DropdownMenuItem>
\end{lstlisting}

\section{Audit Trail \& Immutability}

\subsection{Overview}
Every record in Convex includes:
\begin{itemize}
    \item \texttt{\_id}: Unique identifier (auto-generated, immutable).
    \item \texttt{\_creationTime}: Timestamp of record creation (immutable, set by Convex).
    \item \texttt{author}: Reference to the user who created the record.
\end{itemize}

\subsection{Real-world Translation}
This satisfies \textbf{NIST SP 800-53 AU-2} (Audit and Accountability) and \textbf{ISO/IEC 27001 A.12.4} (Logging and Monitoring).

\subsection{Audit Log Scenario}
An admin investigating a security incident can:
\begin{enumerate}
    \item Query all posts created by a suspicious user.
    \item See the exact \texttt{\_creationTime} of each post.
    \item Cross-reference with access logs (if stored separately).
    \item Determine if a post was modified after creation (by comparing timestamps or using database transaction logs).
\end{enumerate}

\section{Frontend Security: Skeleton Loaders \& UX}

\subsection{Skeleton Loaders}
While loading data from Convex queries, the UI renders placeholder components (``skeletons'') that mimic the final data structure. This prevents the dreaded ``Cumulative Layout Shift'' (CLS).

\subsection{Security Benefit: Preventing Accidental Clicks}
\begin{itemize}
    \item Without skeleton loaders, the page content jumps when data loads, potentially causing a user to accidentally click a ``Delete'' button instead of the intended target.
    \item With skeleton loaders, the layout is stable, reducing ``click hijacking'' vulnerabilities and improving user trust.
\end{itemize}

\subsection{Implementation}
\begin{lstlisting}
// components/custom/elements/admin/usersTable/page.tsx
{users === undefined ? (
  Array.from({ length: 5 }).map((_, i) => (
    <TableRow key={i}>
      <TableCell><Skeleton className="h-4 w-24" /></TableCell>
      <TableCell><Skeleton className="h-6 w-20 rounded-full" /></TableCell>
      <TableCell><Skeleton className="h-4 w-40" /></TableCell>
      {/* ... more skeleton cells ... */}
    </TableRow>
  ))
) : (
  // Render actual data
)}
\end{lstlisting}

\section{SEO, Metadata, \& Trust}

\subsection{OpenGraph (OG) Tags}
The application includes OG meta tags, allowing social platforms (Twitter, Facebook, LinkedIn) to display a preview when users share links.

\subsection{Real-world Benefit}
Explicitly defining OG images and descriptions prevents:
\begin{itemize}
    \item \textbf{Link Preview Manipulation}: An attacker cannot inject a custom image or misleading title into the social preview.
    \item \textbf{Phishing via Social Share}: A legitimate link to \texttt{it-experiment-1.vercel.app} will always show the correct preview.
\end{itemize}

\subsection{Implementation}
\begin{lstlisting}
// app/layout.tsx
export const metadata: Metadata = {
  title: { default: "IT Experiment - 1 | Modern Blog Platform" },
  openGraph: {
    url: "https://it-experiment-1.vercel.app",
    siteName: "IT Experiment - 1",
    images: [{ url: "/og-image.png", width: 1200, height: 630 }],
  },
};
\end{lstlisting}

\section{Infrastructure Security: Vercel Deployment}

\subsection{HTTPS \& TLS}
Vercel automatically provisions and renews TLS certificates, enforcing HTTPS on all connections. Requests to \texttt{http://} are redirected to \texttt{https://}.

\subsection{Real-world Protection}
\begin{itemize}
    \item \textbf{Confidentiality}: All traffic is encrypted with AES-256 or AES-128, preventing \textbf{man-in-the-middle (MITM)} attacks.
    \item \textbf{Integrity}: TLS includes authentication, ensuring the server is authentic.
    \item \textbf{Forward Secrecy}: Modern TLS 1.3 uses ephemeral keys, meaning even if a long-term key is compromised, past sessions remain secure.
\end{itemize}

\section{Comparison to Security Standards}

\begin{table}[h!]
\centering
\begin{tabular}{@{}p{3cm}p{4cm}p{6cm}@{}}
\toprule
\textbf{Standard} & \textbf{Control} & \textbf{Project Implementation} \\ \midrule
\textbf{OWASP Top 10 (2021)} & A01: Broken Access Control & RBAC with client and server-side role checks. \\
& A03: Injection & Zod schemas + Convex type validation. \\
& A06: Vulnerable Components & Using Clerk (third-party auth expert). \\
& A07: Authentication Failures & JWT tokens with secure HTTP-only cookies. \\ \midrule
\textbf{NIST SP 800-53} & AC-1 & Access Control enforced by Clerk and Convex. \\
& AU-2 & Immutable timestamps and author IDs logged. \\
& IA-2 & Clerk's managed MFA and password hashing. \\
& SC-7 & HTTPS enforced by Vercel. \\ \midrule
\textbf{CIS Controls} & Control 1 & Inventory and control (code in Git, dependencies tracked). \\
& Control 5 & Access control via RBAC. \\
& Control 6 & Logging via immutable audit trail. \\ \bottomrule
\end{tabular}
\caption{Detailed mapping of project features to global security standards.}
\label{tab:standards}
\end{table}

\section{Threat Modeling \& Mitigations}

\subsection{Threat 1: Unauthorized Post Creation}
\textbf{Threat:} A regular user creates a post with \texttt{admin} permissions.\\
\textbf{Mitigation:} 
\begin{itemize}
    \item The form validates the \texttt{author} field against the authenticated user's ID (Zod + client).
    \item The server re-verifies that the request comes from the claimed author (Convex mutation handler).
    \item The database enforces a foreign key constraint, ensuring \texttt{author} references a valid user.
\end{itemize}

\subsection{Threat 2: SQL/NoSQL Injection}
\textbf{Threat:} An attacker sends a malformed query to extract data.\\
\textbf{Mitigation:}
\begin{itemize}
    \item Convex abstracts away raw SQL/queries; all queries are type-safe functions.
    \item User input is validated by Zod before submission.
    \item Even if raw queries were used, parameterized queries (used by Convex internally) prevent injection.
\end{itemize}

\subsection{Threat 3: Spam/DoS via Post Flooding}
\textbf{Threat:} A bot creates 1000 posts in a second.\\
\textbf{Mitigation:}
\begin{itemize}
    \item Rate limiting (10-minute cooldown) enforced server-side.
    \item The \texttt{getLastPostTime} query checks the latest post timestamp.
    \item The UI disables the publish button during the cooldown.
\end{itemize}

\subsection{Threat 4: XSS (Cross-Site Scripting)}
\textbf{Threat:} An attacker injects malicious JavaScript in a post description.\\
\textbf{Mitigation:}
\begin{itemize}
    \item React automatically escapes all string interpolations, preventing DOM-based XSS.
    \item If custom HTML were allowed, DOMPurify or similar libraries would sanitize it.
    \item Content Security Policy (CSP) headers can be added to Vercel deployments to prevent inline scripts.
\end{itemize}

\section{Lessons Learned \& Future Improvements}

\subsection{What Worked Well}
\begin{enumerate}
    \item \textbf{Managed Services:} Outsourcing authentication to Clerk eliminated entire classes of vulnerabilities.
    \item \textbf{Type Safety:} TypeScript and Zod caught many errors at compile/runtime, not in production.
    \item \textbf{Real-time Updates:} Convex's real-time subscriptions made the UX feel responsive and secure.
    \item \textbf{Skeleton Loaders:} Improved perceived performance and reduced accidental clicks.
\end{enumerate}

\subsection{Future Security Enhancements}
\begin{enumerate}
    \item \textbf{Server-side RBAC Checks:} Add explicit role validation in every Convex mutation (not just client-side UI checks).
    \item \textbf{Advanced Rate Limiting:} Implement exponential backoff or IP-based throttling for repeated failures.
    \item \textbf{Content Encryption:} Encrypt sensitive fields at rest (e.g., user bios, private posts).
    \item \textbf{Two-Factor Authentication:} Enforce MFA for admin accounts.
    \item \textbf{API Rate Limiting:} Limit Convex query/mutation calls per user per minute.
    \item \textbf{Logging \& Monitoring:} Send all access events to a centralized logging service (e.g., Datadog, Sentry).
    \item \textbf{Penetration Testing:} Hire a professional security firm to test the application.
\end{enumerate}

\section{Conclusion}

\textbf{IT Experiment - 1} serves as a comprehensive blueprint for modern, secure web application development. By strategically combining managed services (Clerk for authentication, Convex for backend), runtime validation (Zod), and responsive UX patterns (skeleton loaders), the platform achieves a ``\textbf{Secure by Default}'' posture.

The project demonstrates that security is not an afterthought or a separate layer—it is woven into every architectural decision:
\begin{itemize}
    \item \textbf{Architecture:} Serverless and managed services reduce the attack surface.
    \item \textbf{Validation:} Multi-layer schemas prevent malformed data.
    \item \textbf{Access Control:} RBAC ensures least privilege.
    \item \textbf{Audit Trail:} Immutable timestamps enable accountability.
    \item \textbf{Moderation:} Real-time tools allow rapid incident response.
\end{itemize}

This holistic approach reflects current industry best practices and provides a foundation for scaling the system while maintaining security guardrails.

\newpage
\appendix
\section{File Structure \& Important Locations}

\subsection{Frontend}
\begin{itemize}
    \item \texttt{app/layout.tsx} — Global layout, metadata, authentication provider.
    \item \texttt{app/page.tsx} — Home page (public feed, user posts, new post form).
    \item \texttt{app/admin/page.tsx} — Admin dashboard.
    \item \texttt{app/security/page.tsx} — Security principles documentation.
    \item \texttt{components/custom/elements/blog/newBlogForm.tsx} — Zod-validated form.
    \item \texttt{components/custom/elements/blog/allPostsFeed/AllPostsFeed.tsx} — Public feed with skeleton loaders.
    \item \texttt{components/custom/elements/admin/usersTable/page.tsx} — User management UI.
    \item \texttt{components/custom/elements/admin/blogsTable/page.tsx} — Blog management UI.
\end{itemize}

\subsection{Backend}
\begin{itemize}
    \item \texttt{convex/types/index.ts} — Schema definitions for users and blogs.
    \item \texttt{convex/functions/query.ts} — Read-only queries (authentication, fetching posts, rate limit checks).
    \item \texttt{convex/functions/mutations.ts} — Write operations (creating/updating posts and users).
\end{itemize}

\subsection{Configuration}
\begin{itemize}
    \item \texttt{public/og-image.png} — Social share preview image.
    \item \texttt{public/favicon.ico} — Branded favicon.
\end{itemize}

\section{References}
\begin{enumerate}
    \item OWASP Top 10 2021 — \url{https://owasp.org/Top10/}
    \item NIST SP 800-53 — \url{https://csrc.nist.gov/publications/detail/sp/800-53/rev-5/final}
    \item CIS Controls — \url{https://www.cisecurity.org/cis-controls/}
    \item Next.js Security — \url{https://nextjs.org/docs/deployment/production-checklist}
    \item Clerk Documentation — \url{https://clerk.com/docs}
    \item Convex Documentation — \url{https://docs.convex.dev}
    \item Zod Documentation — \url{https://zod.dev}
\end{enumerate}

\end{document}

